\subsection{The normalized translation-unit identity}
\label{subsec:ftb-translation-unit-identity}

For the same formal biextension let
\[
 H(X)=h(X,[w_0]X),\quad
 \Theta(X,Y,Z)=
 \exp\bigl(H(X+Y+Z)-H(X+Y)-H(X+Z)-H(Y+Z)
                         +H(X)+H(Y)+H(Z)\bigr),
\]
and
\(\iota(X)=\exp(H(-X)-H(X))\).
The sums and negatives here are those of the formal group.
The canonical-gauge/cubical identities make \(\Theta\) and \(\iota\)
integral units.  Their constant terms are one.

A normalized symmetric translation by a small \(k\)-torsion point
\(\alpha\), at odd line power \(k\), is represented by an integral
unit \(u_\alpha(X)\).
The exact three identities it obeys are
\begin{align}
 u_\alpha(X+Y)u_\alpha(0)
   &=u_\alpha(X)u_\alpha(Y)\Theta(X,Y,\alpha)^k,
                                                        \label{eq:ftb-unit-cube}\\
 (u_\alpha\tau_\alpha^*)^k(1)&=1,             \label{eq:ftb-unit-order}\\
 u_\alpha(0)\iota(\alpha)^k u_\alpha(-\alpha)&=1.
                                                        \label{eq:ftb-unit-sym}
\end{align}
These are \texttt{TrUnit.mul}, \texttt{TrUnit.order}, and
\texttt{TrUnit.sym} in the frozen source.

Here the unit is obtained from the actual integral theta group, not
by extending a generic phase.
The preimage of the integral point in the theta group is a
\(\mathbb G_m\)-torsor on \(\operatorname{Spec}V\), hence has a lift
because \(V\) is local.
There is a unique odd-order symmetric normalization.
For example, if \(k=2r+1\), an arbitrary lift \(w\) has central units
\(a=w^k\) and \(b=w\,\delta(w)\), where \(\delta\) is inversion on
the symmetric theta group; \(a^2=b^k\).
The central rescaling \(t=a b^{-(r+1)}\) satisfies
\(t^k=a^{-1}\), \(t^2=b^{-1}\).
Thus \(tw\) has the required order and symmetry.
The same two equations show uniqueness.

Read this integral lift in the formal line-bundle frame.
Its action is \(f\mapsto u_\alpha\tau_\alpha^*f\), with
\(u_\alpha\in V[[X]]^\times\).
The cubical multiplication identity gives
\eqref{eq:ftb-unit-cube}; its order gives
\eqref{eq:ftb-unit-order}; inversion gives
\eqref{eq:ftb-unit-sym}.
The source realizes the formal frame as the compatible limit of the
actual jets, so the multiplier and the action on regular germs use
the same frame.
Its model comparison is
\texttt{ExistsTranslateUnitSymmetricFormalGroupPoint.model\_unit}
and \texttt{exists\_trUnit}.
The required integral theta-group and square-model comparisons belong
to the actual cubical model, not to a pairing specified only on the
generic fibre.

\begin{srlemma}[The product translation unit is exactly the
\(\kappa\)-ratio]\label{lem:ftb-unit-kappa-ratio}
Suppose \(\mu^*\mu=k1\), \(k\) is odd, \([\mu]\alpha=0\), and
\(\kappa=\exp(S([\mu]X)-kS(X))\) is as in
Lemma~\ref{lem:ftb-kappa}.
For \(n\) normalized translation units set
\(U(X)=\prod_bu_{\alpha_b}(X_b)\).
Then
\[
                            U\,\tau_\alpha^*\kappa=\kappa .
\]
\end{srlemma}
\begin{proof}
All series below belong to the integral domain \(R=V[[X]]\).
Let \(\psi=\tau_\alpha^*\) and let \(E_0,E_\alpha,E_{-\alpha}\)
be evaluation at \(0,\alpha,-\alpha\).
Their values are constants and are fixed by \(\psi\).
Put
\[
 T=\prod_b\Theta(X_b,\alpha_b,\alpha_b)^k,\qquad
 I=\prod_b\iota(\alpha_b)^k .
\]
The definitions of \(\Theta,\iota,\kappa\), and
\([\mu]\alpha=0\), give
\[
 \psi^2\kappa\cdot\kappa\cdot(E_\alpha\kappa)^2T
       =(\psi\kappa)^2E_\alpha(\psi\kappa),\qquad
 E_{-\alpha}\kappa\cdot I=E_\alpha\kappa .
\]
The first is the third-difference identity for \(S\) at
\((X,\alpha,\alpha)\); all terms involving
\(S([\mu]\alpha)\) vanish.
The second is its inversion identity.
Also \(E_0\kappa=1\) and \(\psi^k\kappa=\kappa\).

Equations \eqref{eq:ftb-unit-cube}--\eqref{eq:ftb-unit-sym}, multiplied
over the copies, say
\[
 \psi U\cdot E_0U=U\cdot E_\alpha U\cdot T,\quad
 \prod_{i=0}^{k-1}\psi^iU=1,\quad
 E_0U\cdot I\cdot E_{-\alpha}U=1 .
\]
Define the unit \(\chi\) by
\(\chi\kappa=U\psi\kappa\).
Substitution and cancellation in these identities give
\[
                 \psi\chi\cdot E_0\chi
                                  =\chi\cdot E_\alpha\chi,\qquad
                  \prod_{i=0}^{k-1}\psi^i\chi=1.
\]
The first identity says
\(\psi\chi=(E_\alpha\chi/E_0\chi)\chi\), with constant unit ratio.
Hence the second has the form \(l\chi^k=1\), with \(l\) a constant
unit.  Applying \(E_0\) gives \(\chi^k=(E_0\chi)^k\).
Factor this difference of powers.
The second factor has nonzero constant term
\(k(E_0\chi)^{k-1}\); characteristic zero and the domain condition
therefore imply \(\chi=E_0\chi\).
In particular \(\psi\chi=\chi\).
The inversion and symmetry identities then give \(\chi^2=1\).
The product identity gives \(\chi^k=1\).
As \(k\) is odd, \(\chi=1\), which is the required identity.
This is the full cancellation mechanism of
\texttt{ProdSymmetricTranslationUnitsEqKappaRatio.core};
no root of unity is discarded.
\end{proof}

\subsection{The half-degree bound for the harmonic connection}
\label{subsec:ftb-half-degree-connection}

We record the formal calculation which produces the factor \(1/2\)
in the discriminant loss.
The integral formal biextension has an integral first-order module
\(\mathcal H_V\) with its alternating pairing, an
\(\mathcal O_E\)-action, and an exact sequence
\[
              0\longrightarrow\operatorname{Lie}(\mathcal F)^\vee
               \longrightarrow\mathcal H_V
               \longrightarrow\operatorname{Lie}(\mathcal F)
               \longrightarrow0 .
\]
The action is adjoint for complex conjugation.
In a formal frame a member lifting the tangent direction \(u\) has
the form \((u,\phi-Su)\), where \(u,\phi\) are integral and
\(S=\operatorname{sg}_1(h)\) is the gauge's linear correction.
This describes the source's \texttt{Hset}; the pairing is the
usual Heisenberg pairing
\(\beta((u,v),(u',v'))=v'\cdot u-v\cdot u'\).
The actual chart has an integral basis of tangent eigendirections
\(u_{\tau,i}\).
These structures and their compatibility with the canonical
connection are those of its integral formal biextension.

\begin{srlemma}[Ramified eigenlifts, with paired prime support]
\label{lem:ftb-ramified-eigenlifts}
Assume \(E/\mathbb Q\) is Galois.
Let \(p\) be the chart residue characteristic, write
\(\Delta_E=p^a m\), \(p\nmid m\), and suppose every prime
\(\mathfrak p\mid p\) is distinct from its conjugate.
For every tangent eigendirection there is a decomposition of a full
eigenlift
\[
 (u_{\tau,i},\,\phi_{\tau,i}+\psi_{\tau,i}-Su_{\tau,i}),
 \qquad\phi_{\tau,i}\text{ integral},\qquad
 v(\psi_{\tau,i,j})\ge-\frac{a}{[E:\mathbb Q]}v(p).
\]
There is a partition of the primes above \(p\) into conjugate pairs,
with a choice of side in each pair, such that all pairings between
the integral lifts and errors vanish unless their prime colours are
equal and their sides opposite.
The full eigenlift has zero pairing with a tangent eigendirection
except at the conjugate embedding.
\end{srlemma}
\begin{proof}
Every embedding \(E\to F\) sends \(\mathcal O_E\) into \(V\).
Its residue kernel is a prime \(\mathfrak p\mid p\).
The finite-algebra product decomposition after completing
\(\mathcal O_E\otimes V\) gives integral projectors onto these prime
factors.  Algebraically the same projectors are obtained by Chinese
remainders on the finite algebra; their idempotents lift over the
complete local ring.
Conjugation pairs the factors by the hypothesis
\(\mathfrak p\ne\bar{\mathfrak p}\).
The adjoint action makes two factor summands orthogonal unless their
primes are conjugate.

Choose any integral lift of \(u_{\tau,i}\) in \(\mathcal H_V\), and
apply its integral prime projector.
It still lifts the same tangent direction.
Over \(F\), apply the embedding projector \(e_\tau\).
If \(\omega_j,\omega_j^\vee\) are an integral basis and its trace-dual
basis, this projector is
\[
                 e_\tau=\sum_j\tau(\omega_j)\,\omega_j^\vee
\]
acting on the module.
This identity follows from trace duality, so its sum over all
embeddings is \(1\) and its action is the required eigenprojection.
Choose \(\delta\) in the different with minimal valuation at this
prime.  Then each \(\delta\omega_j^\vee\) is integral, and the
coefficients of \(e_\tau\) have denominators dividing \(\tau(\delta)\).
Since the tangent direction was already an eigenvector, the
difference between the projected and unprojected lifts is vertical,
of the form \((0,\psi)\).
Thus \(\tau(\delta)\psi\) is integral.

The Galois different exponents at primes above \(p\) are equal.
If their ramification and residue degrees are \(e_p,f_p\), and their
number is \(g\), the discriminant identity is
\(a=gf_p d_p\), while \([E:\mathbb Q]=gf_pe_p\).
It follows that
\[
                  v(\tau(\delta))
                  =d_p\,v(p)/e_p
                  =a\,v(p)/[E:\mathbb Q].
\]
Using the weaker inequality is sufficient if a different element is
chosen only to attain the minimum of its image ideal.
This proves the denominator estimate.
Prime-projector orthogonality proves the three required vanishings
for the integral--integral, error--integral and integral--error
pairings.  Eigenprojection and adjointness give the final
embedding-level vanishing.
These are the trace-dual and prime-projector operations of
\texttt{ExistsRamifiedEigenLiftsOfChartDirections} and
\texttt{HolomorphicFamilies/EigRamLift.eigLiftRam\_core}.
\end{proof}

Here is the precise contraction count.
Give every embedding its conjugate-prime colour and its side.
For a monomial \(m\) in one column, put
\[
 \operatorname{pairMin}(m)=
  \sum_{\text{colours }c}
   \min\left\{\sum_{\ell\text{ on side }+,\ c(\ell)=c}m_\ell,
              \sum_{\ell\text{ on side }-,\ c(\ell)=c}m_\ell\right\}.
\]
For a monomial in \(n\) columns, sum this number over the columns.
Twice this sum is at most the total monomial degree, since
\(2\min(A,B)\le A+B\) in each colour and column.
For a polynomial of row degrees \(a\), the total degree is \(d_a\).
Thus its sum of pair minima is at most \(\lfloor d_a/2\rfloor\).

\begin{srlemma}[The paired connection estimate]\label{lem:ftb-paired-connection}
Under the eigenlift and orthogonality conditions of
Lemma~\ref{lem:ftb-ramified-eigenlifts}, let
\(J=\{x\in F:v(x)\ge-c\}\), where
\(c=a\,v(p)/[E:\mathbb Q]\).
For every integral input power series and every integral harmonic
polynomial \(P\) of row degrees \(a_{\tau,i}\), its normalized
harmonic connection value at any line power belongs to
\(J^{\lfloor d_a/2\rfloor}\).
Consequently its valuation is at least \(-c\,d_a/2\).
\end{srlemma}
\begin{proof}
We give the actual connection calculation, including the divided
contractions.  Write \(x=(x_i)\) for the formal-group coordinates and
\(z=(z_i)\) for the auxiliary coordinates, and write \(D_u\) for the
invariant derivative in direction \(u\).  Put
\[
 s_i(x)=\operatorname{sgOf}(h)_i,\qquad
 F_u(x,z)=\sum_i(z_i-s_i(x))u_i .
\]
On the doubled power-series ring the source connection is
\begin{equation}\label{eq:ftb-explicit-connection}
 N_{u,v}=D_u+\sum_i(D_us_i+v_i)\partial_{z_i}-mF_u .
\end{equation}
In particular \(N_{0,\psi}=\sum_i\psi_i\partial_{z_i}\);
it kills a series depending only on \(x\).
A direct commutator calculation gives
\[
 [N_\xi,N_\eta]=m\beta(\xi,\eta).
\]
Indeed invariant derivatives commute,
\(D_uD_{u'}s_i=D_{u'}D_us_i\), and
\(V_{u,v}(F_{u'})=\sum_iv_iu_i'\) for the vector-field part \(V\)
of \(N\).  These identities prove both the commutator and the
vertical killing assertion without a coordinate change.

Let \(H=h(x,[w_0]x)\) and \(E=\exp(mH)\), with \(E(0)=1\).
For an integral lift
\(\xi_h=(u,\phi-Su)\), the canonical-gauge identities say
\[
 D_uH+\sum_is_i u_i\in V[[x]],\qquad
 D_us_i-\sum_jS_{ij}u_j\in V[[x]] .
\]
They follow by differentiating the additive, symmetric and adjoint
gauge-defect identities; \(w_0+\bar w_0=1\) cancels the two linear
adjoint terms.  Consequently the displayed formula gives
\[
 E^{-1}N_{\xi_h}E
  =D_u+\sum_i(D_us_i+\phi_i-(Su)_i)\partial_{z_i}
      +m\bigl(D_uH+\sum_is_i u_i-\sum_iz_i u_i\bigr),
\]
an operator with integral coefficients.
Thus every word in the operators \(N_{\xi_h}\), applied to
\(E f(x)\) with \(f\) integral, is \(E\) times an integral series;
its constant coefficient is integral.
This is the source's \texttt{ConnectionIntegrality} calculation.

Now write the full lift as \(\xi=\xi_h+(0,\psi)\).
Order the two sides with all \(+\) operators to the left of all
\(-\) operators.  Within a side the required operators commute by
the pairing vanishings.  Move every vertical error operator through
the word to the right, where it kills \(E f(x)\).
Each commutation contributes the scalar
\[
 \kappa_{\ell\ell'}
   =\begin{cases}
       m\beta((0,\psi_\ell),\xi_{h,\ell'}),
              &c(\ell)=c(\ell'),\
                 s(\ell)=+,\ s(\ell')=-,\\
       0,&\text{otherwise}.
     \end{cases}
\]
It belongs to \(J\): the tangent component of \(\xi_h\) is integral
and every coordinate of \(\psi\) lies in \(J\).
The resulting ordered Wick formula replaces a polynomial \(Q\) by
\[
       \sum_{j=0}^{\deg Q}
          \frac{\delta_\kappa^j}{j!}Q,\qquad
 \delta_\kappa=\sum_{\ell,\ell'}\kappa_{\ell\ell'}
                         \partial_\ell\partial_{\ell'} .
\]
This formula follows by induction on a word: commuting the next
vertical operator either moves it to the right or selects one
opposite-side letter for a scalar contraction.
The choices of \(j\) contractions are exactly the terms in
\(\delta_\kappa^j/j!\).

This expression introduces no nonintegral coefficient.
For a monomial \(X^AY^B\), indexed on the two sides, its term with
a nonnegative matrix \(N=(n_{ij})\), \(\sum n_{ij}=j\), has scalar
coefficient
\[
 \frac{\prod_i A_i!/(A_i-\sum_jn_{ij})!\,
       \prod_j B_j!/(B_j-\sum_in_{ij})!}
      {\prod_{ij}n_{ij}!}\ \prod_{ij}\kappa_{ij}^{n_{ij}} .
\]
The integer before the product is integral: divide the first
falling-factorial product by \(\prod_{ij}n_{ij}!\) row by row,
obtaining multinomial integers, and leave the second product
unchanged.  Hence every coefficient lies in \(J^j\).
Each contraction removes one letter from either side of one colour.
It reduces \(\operatorname{pairMin}\) by one.
It follows inductively that \(\delta_\kappa^jQ=0\) once \(j\)
exceeds the maximum pair minimum of a supported monomial.
The integral-operator assertion above therefore places the
constant value in the indicated power of \(J\).
Because \(1\in J\), a term in \(J^j\), \(j\le e\), also belongs
to \(J^e\).

It remains to justify that the ordered values compute the desired
harmonic derivative.  Apply the same commutation formula to the
full lifts, splitting them instead as \((u,0)+(0,v)\).
The corresponding contraction coefficient vanishes unless the
embedding labels are conjugate.
For \(n\) columns its second-order operator is a sum over columns
of scalar multiples of
\[
   \sum_{b=1}^n
        \partial_{(\tau,i,b)}
        \partial_{(\bar\tau,j,b)} .
\]
Each such operator kills \(P\) by its harmonicity.
Their sum \(\delta\) therefore kills \(P\), and the finite
exponential \(\sum_j\delta^j/j!\) applied to \(P\) is just \(P\).
The columnwise ordered Wick expansions consequently have exactly
the same total value as the original harmonic connection.
Formally this uses that the second-order operators in distinct
columns commute; expand the finite exponential of their sum and
collect the multinomial coefficients.
In the horizontal case \(v=0\), substitution
\(\operatorname{ev}_s:z_i\mapsto s_i(x)\) obeys the exact chain rule
\[
               \operatorname{ev}_s N_{u,0}
                             =D_u\operatorname{ev}_s .
\]
The \(\sum_iD_us_i\,\partial_{z_i}\) term is precisely the
derivative of this substitution, and \(F_u(x,s(x))=0\).
For a general \((u,v)\), the additional vertical derivative
\(\sum_iv_i\partial_{z_i}\) remains; it was removed only by the
preceding Wick and harmonic comparison.
Since \(s_i(0)=0\), this horizontal substitution preserves constant
coefficients, proving that the resulting constant value is the
original invariant harmonic derivative.  This proves the value comparison as well
as the error estimate.

The argument is precisely the chain
\texttt{ConnectionKillsSeries},
\texttt{ConnectionSecondOfCommute},
\texttt{ExpConjConnectionIntegral},
\texttt{SCWick}, \texttt{SCMain},
\texttt{LemmaHCols}, \texttt{HarmSecond}, and
\texttt{ChartValJ}.
The bound \(2\sum_b\operatorname{pairMin}\le d_a\) now gives
membership in \(J^{\lfloor d_a/2\rfloor}\).
Since \(v(J^s)\ge-sc\), it gives the asserted valuation.
\end{proof}

At a place where \(\Delta_E\) is a unit, the trace-dual projectors are
integral; the same argument has \(J=V\), and the connection value is
integral.
For an actual principal host the arithmetic host data gives
\(E/\mathbb Q\) Galois, discriminant support in
\(\operatorname{cond}E\), and
\(\mathfrak p\ne\bar{\mathfrak p}\) over each prime in that support.
The precise host assertion is
\texttt{HostFieldGaloisDegreeDiscrConjMovesPrimes.host\_galois\_discr\_moves};
its last clause is equivalent to the existence of
\(x\in\mathfrak p\) with \(\bar x\notin\mathfrak p\).
Applying the ramified estimate, or the unit-discriminant estimate,
gives for every line power the uniform chart loss
\begin{equation}\label{eq:ftb-host-chart-loss}
 \rho=\frac{d_a}{2}
  \sum_{p\mid\operatorname{cond}E}
       v(p)\frac{\operatorname{ord}_p\Delta_E}{[E:\mathbb Q]} .
\end{equation}
Terms away from the residue characteristic have valuation zero.
This is \texttt{HarmonicValuationLowerBoundHostFieldAllPlaces.chBound\_host}.
It is the chart bound required by
Proposition~\ref{prop:ftb-mixed-value-bound}.

\paragraph{Source ledger and exact geometric input.}
The translation-unit construction uses the actual formal line frame
and integral theta group on the chart model.
Its producer also uses
\texttt{FormalGroupCharts/SurjectiveAppTopIntegralModelSquare.lean}
to identify the generic and integral-model fibre actions.
This comparison must be supplied by the same N3/model construction.
The pure \(\kappa\)-ratio proof is
\texttt{RingsOfIntegers/ProdSymmetricTranslationUnitsEqKappaRatio.lean}
(blob \texttt{a17745b704f0f016723223e43b95b2f5a284daf5}).
The ramified projection calculation is
\texttt{RingsOfIntegers/ExistsRamifiedEigenLiftsOfChartDirections.lean}
(blob \texttt{dace979b8644da935cf648e5eef6d90c6ab6688d}),
with \texttt{HolomorphicFamilies/EigRamLift.lean}
(blob \texttt{532b04bcef9d8cdeb821bc4b7e792eaa9f9e75fd}).
The formal normal-ordering identity is
\texttt{src/ReciprocityTheory/ReductiveGroups/ArithmeticGeometry/PullbackSections/HolomorphicFamilies/ChartValJ.lean}.
Its source imports the explicit \texttt{WickVanishing.Closed} and
\texttt{WickVanishing.HarmSecond} calculations.
All actual-chart applications require the same integral formal
biextension and its gauge/eigendirection comparison, supplied by
\texttt{JetHarmonics/IntegralChartFormalBiext.lean} and
\texttt{IntegralChartDirectionsAreEigenvectors.lean}.
These are named model joins, not extra numerical assumptions.
